%======================================================================
% CAPÍTULO 03: TARJETA DE IDENTIFICACIÓN MILITAR (TIM)
%======================================================================
% Emisión y gestión de la TIM
% Fundamento Legal: Art. 84 y 90 - Red Sanitaria Militar
%======================================================================

\chapter{Tarjeta de Identificación Militar (TIM)}
\label{ch:tarjeta-tim}

\section{Introducción}

La \textbf{Tarjeta de Identificación Militar (TIM)} es el documento oficial que acredita la condición de miembro de la FANB y es esencial para el acceso a la Red Sanitaria Militar.

\begin{tcolorbox}[
    colback=sandraSoft,
    colframe=sandraDark,
    width=0.9\textwidth,
    halign=center
]
\small
\textbf{Fundamento Legal (Art. 84 y 90)}: La TIM es el credencial material generado por el sistema que permite al efectivo y a sus familiares el acceso directo a la Red Sanitaria Militar. No portarla o no emitirla sistemáticamente imposibilitaría la atención en hospitales militares.
\end{tcolorbox}

\begin{figure}[H]
\centering
\includegraphics[width=0.85\textwidth]{4-tim.png}
\caption{Módulo Principal de Identificación Militar TIM}
\label{fig:modulo-tim}
\end{figure}

\section{Indicadores del Módulo}

El panel de indicadores muestra estadísticas y el termómetro demográfico del afiliado:

\begin{figure}[H]
\centering
\includegraphics[width=0.85\textwidth]{4.1-indicadores.png}
\caption{Panel Analítico de Indicadores de Salud}
\label{fig:indicadores}
\end{figure}

Los indicadores monitorean en tiempo real la situación del afiliado:
\begin{itemize}[label=\textcolor{sandraGreen}{\faIcon{check-circle}}, leftmargin=*]
    \item \textbf{\faIcon{piggy-bank} Cumplimiento de Cotizaciones:} Estado financiero de las cotizaciones obligatorias requeridas para sostener y activar los beneficios fiduciarios.
    \item \textbf{\faIcon{id-card-alt} Régimen de Vigencia:} Ciclo de vida y caducidad oficial del plástico o carnet digital TIM.
    \item \textbf{\faIcon{notes-medical} Historial Semántico de Atención:} Trazabilidad contable e inventario en la Red Sanitaria Militar.
    \item \textbf{\faIcon{users} Beneficiarios Familiares Activos:} Proporción y vigencia técnica del núcleo familiar amparado bajo el titular.
\end{itemize}

\section{Liberación de TIM}

\begin{tcolorbox}[
    colback=white,
    colframe=sandraDark,
    title=\textbf{\faIcon{unlock-alt} Mecanismo de Liberación de TIM},
    width=\textwidth
]
La \keyconcept{Liberación de Identificación} es el protocolo auditable donde el Estado Mayor o el operador levantan la restricción sistémica o física de un carnet inhabilitado.
Se requiere registrar oficialmente el motivo desencadenante del suceso (e.g., Extravío, Robo o Deterioro Severo) para reactivar administrativamente el código individual de cara a una reimpresión en la Red Sanitaria.
\end{tcolorbox}

\begin{figure}[H]
\centering
\includegraphics[width=0.85\textwidth]{4.2-liberacion-tim.png}
\caption{Procedimiento de Liberación del TIM}
\label{fig:liberacion}
\end{figure}

\section{Aprobación de TIM}

\begin{tcolorbox}[
    colback=white,
    colframe=sandraGreen,
    title=\textbf{\faIcon{thumbs-up} Tracking de Aprobación de Carnetización},
    width=\textwidth
]
El ciclo culmina operativamente cuando las autoridades procedimentan el \keyconcept{Seguimiento de Aprobación}. A través de este módulo se dicta el \emph{Status} validado del nuevo carnet, permitiendo el ingreso automatizado y habilitación transaccional expedita en las farmacias y pabellones de toda la República.
\end{tcolorbox}

\begin{figure}[H]
\centering
\includegraphics[width=0.85\textwidth]{4.3-aprobacion-de-tim.png}
\caption{Aprobación Operativa de Identidad}
\label{fig:aprobacion}
\end{figure}

\section{Tipos de TIM}

\begin{longtable}{p{3cm}p{4cm}p{5cm}}
\toprule
\textbf{Tipo} & \textbf{Validez} & \textbf{Características} \\
\midrule
\endhead
\bottomrule
\caption{Tipos de TIM}
\label{tab:tipos-tim}
\endlastfoot

TIM Activa & 5 años & Personal en servicio activo \\
TIM Reserva & 5 años & Personal en la reserva \\
TIM Retiro & Permanente & Personal pensionado \\
TIM Beneficiario & Según vigencia & Familiares dependientes \\
\end{longtable}

\section{Datos en la TIM}

\subsection*{Frente de la Tarjeta}
\begin{itemize}[leftmargin=*]
    \item \textbf{Fotografía}: Imagen del afiliado
    \item \textbf{Nombres y Apellidos}
    \item \textbf{Cédula de Identidad}
    \item \textbf{Grado y Componente}
    \item \textbf{Número de TIM}
    \item \textbf{Código QR}: Para verificación digital
\end{itemize}

\subsection*{Reverso de la Tarjeta}
\begin{itemize}[leftmargin=*]
    \item \textbf{Fecha de Emisión}
    \item \textbf{Fecha de Vencimiento}
    \item \textbf{Estado}: Activo, Suspendido, Vencido
    \item \textbf{Firma Digital}
    \item \textbf{Grupo Sanguíneo}
\end{itemize}

\section{Especificaciones de Fotografía}

\begin{longtable}{p{4cm}p{8cm}}
\toprule
\textbf{Parámetro} & \textbf{Especificación} \\
\midrule
\endhead
\bottomrule
\caption{Especificaciones de fotografía}
\label{tab:foto}
\endlastfoot

\textbf{Formato} & JPG o PNG \\
\textbf{Tamaño} & 10 KB - 500 KB \\
\textbf{Dimensiones} & 354 x 472 píxeles (ratio 3:4) \\
\textbf{Fondo} & Blanco o azul claro uniforme \\
\textbf{Posición} & Rostro frontal, ojos abiertos \\
\textbf{Iluminación} & Uniforme, sin sombras \\
\end{longtable}

\section{Estados de Solicitud}

\begin{longtable}{p{3cm}p{9cm}}
\toprule
\textbf{Estado} & \textbf{Descripción} \\
\midrule
\endhead
\bottomrule
\caption{Estados de solicitud}
\label{tab:estados}
\endlastfoot

Borrador & Solicitud iniciada, incompleta \\
Enviada & Recibida por el sistema \\
En Proceso & En validación de datos \\
Aprobada & En fabricación \\
Despachada & Enviada por correo postal \\
Entregada & Recibida por afiliado \\
Rechazada & Devuelta por errores en datos \\
\end{longtable}

\section{Red Sanitaria Militar}

La TIM habilita el acceso a los servicios de:

\begin{itemize}[leftmargin=*]
    \item \textbf{Hospitales Militares}: Atención hospitalaria completa
    \item \textbf{ Clínicas IPSFANB}: Servicios ambulatorios
    \item \textbf{ Farmacias}: Medicamentos subsidiados
    \item \textbf{ Estudios Médicos}: Laboratorios, imágenes, especialistas
    \item \textbf{Emergencias}: Atención de urgencias 24/7
\end{itemize}

\begin{tcolorbox}[
    colback=sandraSoft,
    colframe=sandraAccent,
    width=0.9\textwidth,
    halign=center
]
\small
\textbf{Importante}: Sin la TIM vigente, el afiliado y sus beneficiarios no pueden acceder a los servicios de la Red Sanitaria Militar. Es responsabilidad del afiliado mantener su TIM actualizada.
\end{tcolorbox}

\newpage
